\chapter{Conclusioni e Sviluppi Futuri}
\label{chap:conclusions}

Il progetto \textbf{DashMe} ha progettato, implementato e validato un ecosistema analitico multi-piattaforma per il monitoraggio integrato della sicurezza alimentare e della conformità regolatoria, connettendo tre grandi basi dati istituzionali mediante Qlik Sense (RASFF), Tableau (FDA) e Microsoft Power BI (EFSA/WHO GEMS). Il contributo principale del lavoro risiede nell'aver dimostrato come l'integrazione di pipeline ETL, modellazione dimensionale e \textit{Visual Analytics} trasformi registri eterogenei e frammentati in indicatori azionabili per il \textit{Project Management} e il \textit{Risk Management} lungo l'intera catena del valore, dalla sorveglianza chimica preventiva in laboratorio fino ai controlli frontalieri e alla logistica inversa dei richiami.

Sulla base delle analisi sviluppate nei Capitoli~\ref{chap:qlik}, \ref{chap:tableau}, \ref{chap:powerbi} e nella sintesi comparativa del Capitolo~\ref{chap:discussion}, è possibile formulare risposte puntuali alle quattro domande di ricerca poste in apertura dell'elaborato:

\begin{itemize}
    \item \textbf{Risposta a RQ1 (Qlik Sense / RASFF -- Geografia dei flussi e barriera doganale europea)}: l'analisi delle 11.787 notifiche valide (2012--2025, con oltre il 98,8\% concentrato dal 2020) dimostra che il rischio alimentare europeo grava per circa l'89\% sul comparto \textit{Alimenti} e presenta una severità prevalente di livello \textit{Grave} (66,8\%, pur in discesa dal 75\%--85\% del 2020--2022 al 50\%--65\% nel 2023--2025). La barriera doganale blocca in media il 10,6\% dei carichi (1.255 respingimenti), lasciando filtrare nei mercati interni quote rilevanti provenienti da esportatori terzi ad alta criticità quali India e Turchia. La topologia Sankey svela una netta bipartizione tra flussi extra-UE intercettati dai grandi gateway portuali di \textit{Germania} e \textit{Paesi Bassi} (1.056 notifiche congiunte) e flussi intra-UE presidiati da auto-notifiche e controlli di prossimità (\textit{Polonia}--\textit{Lituania} e \textit{Francia}--\textit{Belgio}), mentre il rischio biologico è dominato da \textit{Salmonella spp.} (2.893 occorrenze) e dall'hotspot \textit{Pollame e derivati} (1.620 eventi).

    \item \textbf{Risposta a RQ2 (Tableau / FDA -- Geografia industriale, profili clinici, cause/allergeni e governance operativa)}: sui 104.537 prodotti richiamati (37.848 eventi univoci nel periodo 2012--2026) la distribuzione territoriale evidenzia una forte concentrazione in poli produttivi e logistici quali \textit{California} (4.522 eventi), \textit{New Jersey} (2.597) e \textit{Florida} (2.462), mentre a livello settoriale emerge una netta dicotomia tra volumi e letalità: \textit{Dispositivi} (40.311; 88,33\% Classe~II) e \textit{Medicinali} (18.048; 80,76\% Classe~II) dominano i volumi, ma la severità estrema di Classe~I si concentra in \textit{Prodotti per Animali} (47,65\%) e \textit{Cibo e Bevande} (43,57\%). Sulle sei macro-cause epidemiologiche (43.856 prodotti), dopo il picco storico del 2015 (4.846 prodotti), il minimo pandemico del 2021 (2.037) e il rimbalzo del 2023 (3.312), le \textit{adulterazioni farmaci} guidano per numerosità (17.574; 86,49\% Classe~II), mentre \textit{Salmonella} (86,68\% Classe~I) e \textit{Listeria} (82,09\% Classe~I) determinano il rischio letale, dominando anche i regimi della catena del freddo (55,8\% nel \textit{Refrigerato / Fresco} e 52,4\% nel \textit{Surgelato}) e i materiali flessibili in \textit{Plastica / Busta} (8.395 prodotti, contro i 4.800 in \textit{Vetro / Bottiglia / Barattolo} trainati al 74,4\% da adulterazioni farmacologiche); sul fronte allergenico (7.954 prodotti nella mappa dedicata), l'analisi isola i profili principali (\textit{Latte / Latticini} 2.084, \textit{Frutta a guscio} 870, \textit{Soia} 726) pur evidenziando l'artefatto lessicale \texttt{"peanut"}/\texttt{"nut"} che convoglia 1.046 richiami da arachidi in \textit{Più Allergeni Contemporaneamente} (2.766). Infine, la permanenza di 18.555 prodotti in stato ``In corso'' (17,75\%) e la concentrazione di Classe~I nelle reti \textit{Regionale (Multi-Stato)} (32,43\%) e \textit{Internazionale / Export} (20,99\%) individuano il principale nodo critico di governance post-vendita.

    \item \textbf{Risposta a RQ3 (Power BI / EFSA-WHO -- Sorveglianza chimica preventiva e driver di positività)}: su 2.424.475 rilevazioni laboratoristiche (1984--2017), il tasso di positività effettivo si attesta al 39,63\% globale (959.424 determinazioni $>$ LOD su 2.421.145 misure valide, escludendo i 3.330 nulli) e al 43,44\% sul sottoinsieme dei 32 Stati nazionali (1,62Mln analisi su 724,23K campioni). L'analisi contrappone l'elevata capacità volumetrica della \textit{Germania} ($\approx 638\text{K}$ analisi; $\approx 38,1\%$ positività) ai programmi mirati ad alta incidenza dei Paesi nordici ($\approx 68\%$--$85\%$), evidenziando il bioaccumulo di \textit{Cadmio} nelle frattaglie suine e bovine (65.604 positività nel comparto carneo) e la polarizzazione tossicologica tra \textit{Diossine e Furani} (99,99\% in fascia \textit{Ultra-Tracce} 0--1~$\mu\text{g/kg}$) e \textit{Contaminanti di Processo} (38,89\% in fascia \textit{Elevata} e 16,78\% in fascia \textit{Critica}). Sul piano operativo, i pannelli di \textit{Screening Multi-Residuo ($>10$ analiti)} generano circa 540K analisi e 290K positività da soli 20K campioni fisici, mentre l'apprendimento automatico integrato (\textit{Key Influencers}) isola moltiplicatori di probabilità fino a 2,53x e segmenti decisionali con incidenza fino al 72,9\%.

    \item \textbf{Risposta a RQ4 (Confronto multi-piattaforma -- Trade-off tra Qlik Sense, Tableau e Power BI)}: nessuna singola piattaforma massimizza simultaneamente tutte le esigenze della filiera analitica. Qlik Sense eccelle nell'esplorazione associativa dei flussi di rete e delle correlazioni non anticipate (Sankey e analisi dei valori esclusi); Tableau offre la massima flessibilità visiva e narrativa per scomporre gerarchie complesse su eventi testuali; Power BI garantisce la superiore scalabilità colonnare su milioni di righe, una rigorosa disciplina di modellazione dimensionale (\textit{Star Schema} + DAX) e l'integrazione nativa di componenti diagnostici di \textit{machine learning}. L'impiego sinergico dei tre ambienti conferma quindi il valore di un'architettura BI pluralistica e orientata allo specifico dominio del dato.
\end{itemize}

In prospettiva futura, i risultati e i limiti documentati nel Capitolo~\ref{chap:discussion} tracciano tre linee prioritarie di evoluzione del progetto DashMe:
\begin{enumerate}
    \item \textbf{Automazione incrementale dell'ingestione tramite API e portali ufficiali}: sostituire i caricamenti statici basati su snapshot CSV e XLSX con connettori schedulati verso gli endpoint REST di \textit{openFDA}, i feed pubblici del portale \textit{RASFF} e le basi dati periodiche \textit{EFSA / WHO GEMS}, abilitando il monitoraggio quasi in tempo reale dei richiami ``In corso'' e l'aggiornamento incrementale delle serie storiche.
    \item \textbf{Correzione mirata dei difetti di pipeline emersi nell'audit}: intervenire puntualmente sulle anomalie individuate nei tre ambienti, ovvero: (a) in \textit{Tableau}, sostituire la condizione \texttt{"nut"} con una corrispondenza disambiguata (es.\ \texttt{"tree nut"}) per eliminare il convogliamento improprio dei 1.046 richiami da arachidi in \textit{Più Allergeni Contemporaneamente}, assegnare codici colore distinti a \textit{adulterazioni farmaci} e \textit{Salmonella} eliminando la collisione su \texttt{\#529ee0}, e completare il consolidamento anagrafico dei gruppi multi-sito come \textit{Medline Industries LP}; (b) in \textit{Power BI}, filtrare o convertire preliminarmente i record con \texttt{UnitID = 1} per evitare distorsioni sul 95\textdegree{} percentile della Francia e riallineare le etichette brevi (\texttt{Gruppo Alimentare (breve)}) tra modello \texttt{.pbix} e documenti esportati; (c) in \textit{Qlik Sense}, introdurre una pre-normalizzazione dell'escaping dei doppi apici nel campo \texttt{subject} (recuperando le 4 righe scartate), mappare le 139 occorrenze non classificate (\texttt{-}) e allineare la legenda del grafico stagionale alle cinque serie visualizzate.
    \item \textbf{Modellazione predittiva cross-dataset della severità e della positività}: sfruttare le feature ingegnerizzate lungo i tre moduli (origine geografica, matrice alimentare, profilo di packaging, regime termico e ampiezza dei pannelli analitici) per addestrare modelli di classificazione supervisionata capaci di stimare precocemente la probabilità di esito positivo in laboratorio o di classificazione in Classe~I al momento della notifica iniziale, supportando una sorveglianza proattiva e basata sul rischio.
\end{enumerate}
